\documentclass[12pt]{article}
\usepackage[utf8]{inputenc}
\usepackage[T1]{fontenc}
\usepackage{hyperref}
\title{Wichita Language Resources for K-12 Curriculum Development}
\author{Compiled by Claude Sonnet 5.5 (claude-sonnet-5-5) for Andr\'as Kornai}
\date{}
\begin{document}
\maketitle

\section{Introduction}
Wichita (ISO 639-3 wic; Glottocode wich1260) is a Caddoan language of the Wichita and Affiliated Tribes, based in Anadarko, Oklahoma. The last fluent speaker, Doris Jean Lamar-McLemore, died in 2016. A 2018 thesis argues the language is not extinct because semi-speakers and documentation remain. The tribe has launched language programme work, and a heritage and revitalization curriculum is the natural goal.

\section{Primary Educational Resources}
The Language Conservancy's Wichita Vocab Builder is a free app with 26 categories and more than 464 words and phrases, with audio for each word. It was designed with Wichita speakers: \url{https://apps.apple.com/app/id6464505392}.
The Language Conservancy lists the Wichita partnership: \url{https://languageconservancy.org/languages-we-work-with/}.

\section{Linguistic and Reference Materials}
David Rood has documented and described Wichita for decades. Overviews of the grammar are cited as Rood (1976, 1996) and Mirzayan (2008); a source link for these was not found in search results, so none is given here.
The University of Oklahoma holds a thesis on Wichita revitalization (Sapphire Hodges, 2018): \url{https://shareok.org/handle/11244/299837}.
A Wichita-related page from the Max Planck DoBeS archive: \url{https://dobes.mpi.nl/posters/print/americas-wichita-print.pdf}.

\section{Religious and Community Resources}
The Sam Noble Museum at the University of Oklahoma holds the Wichita Language Program Collection (161 items, 1999 to 2014), including a Wichita Audio Dictionary. Access should be confirmed with the museum: \url{https://samnoblemuseum.ou.edu/collections-and-research/native-american-languages/native-american-languages-collections/wichita-language-program-collection}.
The Endangered Languages Project has a page on language learning and teaching: \url{https://endangeredlanguages.com/node/109792}.

\section{Digital and Audio Resources}
The Vocab Builder app (above) includes audio. The EBSCO research starter reports that the tribe launched an audio-enabled online dictionary, but its address was not found in search results: \url{https://ebsco.com/research-starters/social-sciences-and-humanities/wichita-tribe}.

\section{Curriculum Development Context}
There are no first-language speakers, so a curriculum would be a heritage and revitalization course built from archived recordings, the app vocabulary and published grammar. The tribe has reportedly secured funding for language staff, beginner materials and a five-year plan. Any school programme should be developed with the Wichita and Affiliated Tribes, who control access to the language materials.
\end{document}